% One-page A4 summary extracted from presentation speaker notes
\documentclass[a4paper,11pt]{article}
\usepackage[utf8]{inputenc}
\usepackage{geometry}
\usepackage{parskip}
\usepackage{hyperref}
\geometry{margin=1in}

\begin{document}
\begin{center}
  {\LARGE Dynamic Attack Graphs: Incremental Re-computation with Differential Dataflow}\\[6pt]
  {\large Stefania‑Cristina Mozacu; Emil C. Lupu}\\[12pt]
\end{center}

\section*{Executive summary}
This study implements a proof-of-concept system for incremental recomputation of attack graphs using the Timely + Differential Dataflow Rust libraries. Input facts (vulnerabilities, network edges, firewalls, attacker positions, goals) are translated into keyed collections and derived facts are computed using joins, antijoins, and iterative fixed-point computation. The approach reduces update cost for local changes in shallow topologies (star) while worst-case chain-like dependencies remain costly.

\section*{Key components}
\begin{itemize}
\item Data model: \texttt{VulnerabilityRecord}, \texttt{NetworkAccessRule}, \texttt{FirewallRuleRecord}, \texttt{AttackerStartingPosition}, \texttt{AttackerTargetGoal}
\item Derived facts: \texttt{EffectiveNetworkAccess}, \texttt{AttackerCodeExecution}, \texttt{AttackerOwnsMachine}, \texttt{AttackerGoalReached}
\item Operators: \texttt{antijoin} (for denies), \texttt{iterate} (recursion), \texttt{distinct} (deduplication)
\end{itemize}

\section*{Benchmarks (summary)}
Topologies: star, chain, random-cut on chain. Metrics: initial full recomputation time, incremental update time, and speedup. Representative outcome: star topology shows large speedups; chain shows marginal improvements for cuts near the head.

\section*{Reproducibility}
Place the Graphviz images (\texttt{graph_initial.png}, \texttt{graph_final.png}) in the project root's working directory. Key commands:
\begin{verbatim}
cargo build --release
cargo run --release --bin attack-graph
dot -Tpng graph_initial.dot -o graph_initial.png
docker build -t attack-graph .
docker run --rm attack-graph
\end{verbatim}

\section*{Next steps}
Scaling to multiple workers, integrating real vulnerability scanner output, and applying the approach to real network topologies.

\section*{Location of code}
See the repository root. Key files: \texttt{src/rules.rs}, \texttt{src/benchmarks.rs}, \texttt{Dockerfile}, \texttt{graph_initial.dot}

\end{document}
